% 第2章 相关工作

\section{视觉语言模型}

视觉语言模型（Vision-Language Models，VLMs）是连接计算机视觉与自然语言处理的桥梁技术，通过将视觉编码器与大语言模型对齐，实现跨模态图像理解和文本生成能力。近年来，该类模型的快速发展为三维场景的自动化描述提供了技术基础。

在开源视觉语言模型方面，LLaVA（Large Language and Vision Assistant）\cite{llava} 是代表性的工作之一。其核心架构由 Vision Transformer（ViT）视觉编码器和 LLaMA 语言解码器组成，通过一个简单的线性投影层将视觉特征映射至语言空间。LLaVA v1.5 7B 版本在视觉指令微调后，能够对任意图像区域生成自然语言描述。然而，LLaVA 的描述质量高度依赖于输入图像的质量和分辨率。在本文的三维物体裁剪场景中，由于物体裁剪区域往往较小（如仅占原始图像 5\%—15\%），且包含大量背景噪声，LLaVA 生成的描述常常偏向于整体场景氛围而非物体本身属性，这一局限性成为本文后续 prompt 优化的直接动机。

Qwen2.5 系列模型\cite{qwen2_5} 是阿里巴巴 Qwen 团队发布的开源大语言模型家族，涵盖 0.5B 至 72B 参数的多个规模变体。其中 Qwen2.5-3B 版本在推理能力和资源消耗之间取得了优秀的平衡：在 MMLU 基准上达到 65.6 分，同时推理速度远快于 7B 及以上版本。本文在推理引擎选型实验中系统测试了 Qwen2.5 系列的三个变体（3B、7B、14B），验证了小模型在结构化数据合成任务中的充分性。Qwen2.5 的 GGUF 量化格式（Q4\_K\_M）进一步将模型压缩至 1.9—4.4 GB，使得在单张消费级 GPU 上的本地化部署成为可能。

An 等人（ICLR 2025）的研究表明，将文本等额外模态信息与三维点云数据融合，可以显著提升 few-shot 三维语义分割的性能\cite{multimodal_3d}。该工作使用预训练的 CLIP 模型提取文本特征，通过跨模态注意力机制将语言先验知识注入点云分割网络，在 ScanNet 和 S3DIS 数据集上均取得了 state-of-the-art 的 few-shot 性能。这一发现提示了多模态融合对三维理解任务的重要性，也为本文后续引入多模态描述融合提供了理论基础。

\section{思维链推理}

思维链（Chain-of-Thought，CoT）提示是一种通过显式引导语言模型生成中间推理步骤来提升复杂任务性能的方法。其核心思想是模仿人类解决问题的过程——将复杂问题拆解为多个可处理的子步骤，逐步推导出最终答案。在 GSM8K 数学数据集上，PaLM-540B 使用 8 个 CoT 示例后准确率从 34\% 提升至 57\%，超越了当时微调后的 GPT-3 175B，证明了 CoT 在多步推理任务中的有效性。

然而，传统 Few-shot CoT 存在明显局限。首先，人工设计的 CoT 示例不仅耗时耗力，且示例质量直接影响推理效果。Zhang 等人提出的 Auto-CoT\cite{auto_cot} 通过聚类采样自动生成演示：将待处理问题聚类为若干类别，从每个类别中选择代表性样本，使用语言模型为这些样本自动生成推理链。Auto-CoT 在多个推理基准上接近甚至超越了人工设计 CoT 的效果，验证了自动化示例生成的可行性。

Cheng 等人对新一代强大语言模型上的 CoT 提示进行了系统性重新审视\cite{chain_of_thought}。该研究发现了一个重要的反直觉现象：对于 Qwen2.5 等强大语言模型，Zero-Shot CoT 在 GSM8K 和 MATH 数据集上的推理准确率普遍高于传统 Few-shot CoT。进一步的替换实验揭示：传统 CoT 示例的核心作用仅是规范输出格式（如引导模型使用特定格式包裹答案），而非传递推理知识。注意力机制可视化进一步证实：模型在解码时几乎不关注示例区域，而是聚焦于问题指令本身。这一发现为本文采用零样本结构化 prompt 进行场景描述优化提供了直接的理论支撑——对于能力足够强大的语言模型，精心设计的格式约束比提供示例更有效。

Zhou 等人探索了使用大语言模型为软件架构决策生成设计理由的可行性\cite{design_rationale}。其研究结果表明零样本方法的平均精确度（0.278）和 F1 值（0.381）均优于 CoT 方法（0.237 和 0.351），且推理速度最快、资源消耗最低。CoT 生成的论证虽然数量增加 26.7\%，但其中许多论证流于表面、缺乏深度分析，进一步验证了零样本推理在结构化输出任务中的优势。

\section{Code2Logic：游戏逻辑驱动的数据合成}

Code2Logic 由 Tong 等人提出（NeurIPS 2025），是本文方法论的核心来源\cite{code2logic}。该方法创造性地将游戏源代码作为高质量推理数据的自动化生成源，提出了三阶段数据合成框架。

第一阶段为游戏代码构建。使用大语言模型根据游戏规则描述生成可执行的游戏源代码。以推箱子为例，仅需单行自然语言提示，LLM 即可输出包含状态空间定义、合法动作集合、胜负判定条件的完整 Python 代码，建立了从自然语言规则到程序语义的精确映射。

第二阶段为问答模板设计。从游戏逻辑中提取三类渐进式推理任务：（1）Target Perception（目标感知）——聚焦静态元素识别；（2）State Prediction（状态预测）——模拟动作执行后的状态变化；（3）Strategy Optimization（策略优化）——在多步推理中寻求最优解。这三类任务构成认知梯度的渐进式推理挑战。

第三阶段为数据引擎构建。将游戏核心代码作为数据生成引擎，通过自动化参数采样和状态快照，批量生成带视觉轨迹的问答对。单款游戏平均 7.5 小时的初始开发投入后，数据引擎可零边际成本无限生成新样本，单位标注成本仅为人工标注的 1/50。

Code2Logic 构建的 GameQA 数据集涵盖 30 种游戏、158 项认知任务和 14 万问题。实验揭示了几项关键洞察：（1）仅用 GameQA 训练的 VLM 在 7 个主流视觉基准上平均提升 2.33\%；（2）使用 20 款游戏训练的模型比 4 款游戏的域外任务准确率提升 1.8\%，验证了游戏多样性对泛化能力的驱动作用；（3）GRPO 强化学习策略使视觉感知与文本推理错误率分别降低 19.2\% 和 11.5\%。

Code2Logic 的核心局限在于其完全依赖游戏代码的显式逻辑结构——游戏世界中物体的位置和交互规则均由代码明确定义。而真实物理场景中，这些信息需要通过传感器数据和视觉推理间接获取，噪声、遮挡和光照变化都会引入不确定性。本文的工作正是试图桥接这一差距：将模板化数据合成方法论从二维游戏世界迁移至三维真实场景。

\section{三维场景图构建}

三维场景图（3D Scene Graph）将场景表示为结构化图——节点代表物体实例，边代表空间或语义关系。与二维场景图不同，三维场景图需要处理深度信息、几何约束和视角变化等额外挑战。

Concept-Graphs\cite{conceptgraphs} 是三维场景理解领域的重要工作，通过整合多种感知模块实现从 RGB-D 到三维场景图的端到端构建。其三个核心组件分别为：GradSLAM\cite{gradslam} 提供可微分的 SLAM 前端，从 RGB-D 序列和相机轨迹重建三维场景几何；Grounded-SAM\cite{grounded_sam} 组合 Grounding DINO 的开放词汇检测和 SAM 的精细分割，实现任意物体的检测与分割；CF-SLAM 将二维分割结果通过相机参数投影至三维空间，经重叠检测、相似度匹配和 DBSCAN 聚类，构建物体级场景图。

Replica 数据集\cite{replica} 由 Straub 等人发布，包含 18 个高精度三维重建的室内场景。每个场景提供约 2000 帧 RGB 图像（640×480）、对应深度图和由运动捕捉系统采集的相机轨迹。本文使用 room0、room1 和 office0—4 共 7 个场景，覆盖起居室和办公室两类室内环境。

\section{合成数据与域随机化}

合成数据的核心思想是通过计算机图形学自动生成带标注的训练数据。其在计算机视觉领域的核心优势包括：标注自动且完美准确、场景参数完全可控、可无限扩展。

域随机化（Domain Randomization）在合成数据基础上进一步提升模型泛化能力：在仿真环境中对非本质属性（光照、颜色、纹理等）进行随机扰动，使模型学到对这些变化不敏感的鲁棒特征。本文在 Isaac Sim 中应用了三种域随机化：物体位置随机化、材质颜色随机化、光照参数随机化。

NVIDIA Isaac Sim\cite{isaac_sim} 是面向机器人仿真和数据合成的综合性平台，基于 USD（Universal Scene Description）格式，支持 RTX 光线追踪渲染和无头模式 Vulkan 后端。在本文中，Isaac Sim 被用作标注质量的验证工具——以合成场景中的精确 ground truth 为基准，对 Concept-Graphs 管线的推断结果进行可量化的精度评估。

\section{本章小结}

综上所述，视觉语言模型为三维场景理解提供了基础能力，思维链推理研究明确了零样本结构化输出的优越性，Code2Logic 验证了模板化数据合成的工程可行性，三维场景图技术提供了从传感器数据到结构化表示的完整链路，合成数据方法则为标注质量的闭环验证奠定了基础。本文的创新正在于将这些技术线交汇融合——将游戏逻辑驱动的合成范式从二维游戏世界迁移至三维真实场景，构建可验证的自动标注系统。
